\section{相关工作}

\subsection{直方图均衡与 CLAHE}
全局直方图均衡（HE）通过重新分配灰度分布提升整体对比度，但在低照度场景下往往会导致高亮区域过曝、暗区噪声被同时拉升。为缓解这一问题，Pizer 等提出了 AHE 及其变体\cite{pizer1987ahe}，随后 Zuiderveld 将对比度限制机制引入局部直方图均衡，形成了 CLAHE\cite{zuiderveld1994clahe}。这一路线的优点在于无需训练、实现简单、局部对比度提升直接；缺点则在于固定 clip limit 难以同时适应平坦暗区、弱纹理区和高结构区，容易出现噪声放大或块状伪影。本文工作正是在这一传统框架内继续推进，希望通过区域自适应的 clip limit 分配来改善固定参数 CLAHE 的局限。

\subsection{Retinex 与照明估计类低照度增强}
除直方图均衡外，另一类经典思路是通过照明-反射分解或照明估计来恢复低照度图像。Jobson 等提出的多尺度 Retinex（MSR）方法\cite{jobson1997retinex} 通过多尺度滤波增强局部细节，在低照度增强中具有重要启发意义。随后，Wang 等提出 NPE 方法\cite{wang2013loe}，强调增强过程中对自然度的保持；Fu 等提出基于图像融合的弱照度增强方法\cite{fu2016fusion}；Guo 等提出 LIME\cite{guo2017lime}，通过估计照明图实现更稳定的亮度恢复。近年来，Wei 等又将 Retinex 分解与深度学习结合，提出了 Retinex-Net\cite{wei2018retinexnet}。总体而言，这类方法通常能够提供比简单直方图均衡更灵活的亮度恢复机制，但也会面临自然度保持、参数设置复杂以及暗区噪声同步增强等问题。

\subsection{增强与去噪联合方法}
低照度图像的一个核心困难在于：亮度增强和噪声抑制往往相互耦合。若仅追求亮度和对比度提升，噪声也会被同步放大，从而破坏纹理与自然度。围绕这一问题，Li 等提出了同时兼顾去噪与对比度扩展的增强方法\cite{li2015lowlight}；Ren 等进一步从顺序分解的角度联合建模增强与去噪过程\cite{ren2018joint}。这些工作表明，低照度增强不应只看“是否变亮”，还需要关注增强后的暗区稳定性与噪声控制。本文自定义的 NAR（Noise Amplification Ratio）正是基于这一考虑，用于刻画增强前后暗区噪声放大的相对程度。

\subsection{亮度-色度分离与本文定位}
在实现层面，许多传统增强方法会将图像转换到 YCbCr、HSV 等亮度-色度分离空间，仅对亮度通道进行处理，以降低直接在 RGB 空间增强带来的色偏风险。本文同样采用这一思路，仅对亮度通道进行自适应增强，并将色度补偿降级为真实图主观展示中的可选后处理，而非合成数据定量实验中的核心模块，以减少颜色变化对 RGB-SSIM 的额外干扰。与 Retinex 类或深度学习类方法不同，本文不试图提出新的低照度增强范式，而是在 CLAHE 路线内，围绕噪声放大、块效应和 clip limit 固定不变这三个问题，构建一种更可解释、便于消融分析的局部结构感知改进方案。
